\chapter[Learning and persistent uncertainty]{Learning, persistent uncertainty, and who knows what}
\label{sim:learning}

A future disclosure may resolve an uncertainty about today's company. It may also reveal nothing relevant. Simulation must allow both outcomes. It must distinguish a change in the company from a change in the analyst's knowledge, and distinguish the analyst's knowledge from information available to managers, lenders and investors. These distinctions determine which actions are possible and which apparent forecasting successes are merely uses of future information.

\section{Five dates attached to one fact}
A loan can be signed on 15 January, take economic effect on 1 February, appear in a report covering the quarter ending 31 March, and become public on 10 May. A system may ingest the report on 12 May. The signature, economic, reporting, publication and ingestion dates answer different questions. A forecast made on 30 April cannot use the 10 May disclosure simply because it describes a March balance.

For a report-based forecast origin $t$, let $D_t$ contain assertions available by $t$ under the declared availability rule. If the analysis uses public information, publication time is relevant. If it models a lender with private contractual access, the information rule differs. Ingestion time controls reproducibility of the system's own operation; it need not equal the earliest public availability.

A historical replay should therefore specify both the economic date and evidence cutoff. A corrected extraction performed today may accurately recover a number that was publicly available before the old cutoff. A new restatement published today was not available then. The first can repair a data-processing mistake in an old evidence view; the second changes the information set and must not be silently backfilled into the old forecast.

\section{Persistent unknown facts}
Let $\theta$ denote a contract characteristic fixed at inception, such as whether a loan is fixed or floating. Let $X_t$ denote changing economic state, including cash, accrued interest and remaining principal. A simple transition is
\begin{equation}
 X_{t+1}=T(X_t,\theta,a_t,\varepsilon_{t+1}),\qquad \theta_{t+1}=\theta_t.
 \label{sim:persistent-theta}
\end{equation}
Here $a_t$ is a decision and $\varepsilon_{t+1}$ a future innovation under the specified scenario or stochastic process. The second equation is essential. Uncertainty about a fixed contract does not license redrawing it each period. Refinancing or amendment can change the terms, but those are dated events with their own permissions and cash consequences.

Suppose a loan is either fixed at five per cent or floating at the market rate, and the market rate is seven per cent for each of two future years. The possible interest sequences are $(5,5)$ and $(7,7)$. Independent yearly resampling adds $(5,7)$ and $(7,5)$. Those additional paths describe contract changes that the model did not authorize. Their inclusion can alter liquidity minima, covenant duration and refinancing behavior.

The same persistence issue affects customer identity, asset vintage, tax basis, priority and collateral. A compact state must retain persistent variables even if current totals do not depend on them. Eliminating a variable because it is irrelevant today is unsafe if it changes later observations or events.

\section{Set-valued evidence updates}
When no probability measure is justified, retain a feasible family $\mathcal H_t$. If a new compatible hard observation $y$ supplies condition $C_y$, updating the same historical scope gives
\begin{equation}
 \mathcal H_{t}^{+}=\{h\in\mathcal H_t:C_y(h)\}.
 \label{sim:set-update}
\end{equation}
This operation narrows the old family. It does not advance the economic clock. Forward propagation then maps each remaining history into future possibilities under the transition relation. A smaller historical family can still produce a wider future range when the horizon lengthens or new shocks are admitted.

An empty intersection calls for investigation. A source may conflict with another source, the observation mapping may be wrong, or a model assumption may be too narrow. Selecting one convenient history to restore nonemptiness would hide the contradiction. The result should identify the conflicting premises and preserve the previous view.

A restatement is not always an additional constraint on the old asserted amount. It may replace that assertion in a current evidence view. The system should preserve the original source and link the correction to it. Historical replay can then reproduce the original view, while current reconstruction uses the revised one. Conjoining both as exact truths would manufacture infeasibility.

\section{Probabilistic filtering requires a model}
If a prior measure $\mu_t$ and likelihood $L(y\mid h)$ are admitted, Bayesian updating has the form
\begin{equation}
 \mu_t^+(dh)=\frac{L(y\mid h)\mu_t(dh)}{\int L(y\mid u)\mu_t(du)}.
 \label{sim:bayes-update}
\end{equation}
The denominator must be finite and positive. Its role is normalization, not an opportunity to substitute an arbitrary fallback when the observation has zero modeled probability. A zero denominator can reveal a misspecified likelihood, an impossible observation under the model or inappropriate handling of a continuous conditioning event.

For continuous exact observations, ordinary density multiplication must be understood relative to the declared dominating measure and conditional distribution. Adding a tiny Gaussian error can make a computation run, but it changes the observation model. Its scale should reflect a justified measurement or reporting process, or be explicitly treated as a sensitivity assumption. Numerical convenience does not establish the company's measurement error distribution.

The POMDP formulation in \citet[Sections 3.1--3.4]{kaelbling1998} provides a formal treatment when transition and observation probabilities are known. Its belief state summarizes the history needed for control under those assumptions. An unweighted feasible family is not a probabilistic belief state merely because it is stored in the same data structure. Supporting both requires distinct semantics.



\section{A finite filter that preserves contract identity}
Let $\theta$ be an unknown but fixed contract type, $x_t$ a changing hidden
operating state, and $y_t$ the analyst's observation. A declared probabilistic
model supplies transition probabilities $P_\theta(x'\mid x)$ and observation
likelihoods $\ell_t(y\mid\theta,x')$. With joint posterior weights
$w_{t-1}(\theta,x)$, the prediction and update are
\begin{align}
w_t^-(\theta,x')&=\sum_x
  P_\theta(x'\mid x)w_{t-1}(\theta,x),\\
Z_t&=\sum_{\theta,x'}\ell_t(y_t\mid\theta,x')w_t^-(\theta,x'),\\
w_t(\theta,x')&=\ell_t(y_t\mid\theta,x')w_t^-(\theta,x')/Z_t.
\end{align}
There is no summation that changes $\theta$ into another contract type. A
contract modification needs an explicit event and a different transition
model. If $Z_t=0$, the result is an evidence conflict under the declared model;
normalizing arbitrary fallback weights would answer another question.

An uninformative observation has constant likelihood across states, so its
update does not change the predicted relative weights. A missing disclosure
is uninformative only when the observation model makes it so. If no prior is
justified, propagate a feasible set of joint $(\theta,x)$ pairs through the
transition relation and evidence restrictions. That procedure produces support,
not posterior probabilities. A declared set of priors can instead be updated
member by member; zero-likelihood members need an explicit conditioning rule,
and one must not replace the result by independently chosen period marginals.

A source correction can remove an earlier restriction. Retaining only the
last posterior may then be insufficient: a zeroed type cannot be recovered
by multiplying it by another likelihood. Keep the original prior or feasible
family, the evidence versions and the transition model, and replay all
currently active evidence. Checkpoints can accelerate replay only when their
evidence dependencies remain valid. The reference filter compares this
recursion with direct path enumeration for four priors and all 27 length-three
observation sequences. It also checks rare positive types, uninformative
observations, conflict, and a correction that reopens a discarded type. This
is an exact finite method; estimating a defensible issuer likelihood and
controlling particle approximation remain separate problems.


\section{An informative and an uninformative disclosure}
Return to the fixed-versus-floating loan. Assume, only for this example, prior probabilities of one half each. A report can say ``floating'' or ``fixed.'' Suppose the report identifies the correct type with probability 0.9 and the wrong type with probability 0.1, independently of other modeled quantities conditional on the type. After a ``floating'' report, Bayes' rule gives
\begin{equation}
 \Pr(F\mid y=F)=\frac{0.9(0.5)}{0.9(0.5)+0.1(0.5)}=0.9.
\end{equation}
These probabilities are supplied assumptions, not estimates from the loan balance. The example isolates the mechanics of updating. If the report were equally likely under both types, its likelihood would cancel and the prior would remain unchanged. More pages of disclosure need not imply more information about the exposure of interest.

In the set-valued version, an exact contract note may eliminate one type. An ambiguous phrase compatible with both eliminates neither. A probabilistic false-positive model is a different treatment from a set of possible interpretations. The analyst should choose the treatment supported by the evidence process and preserve its assumptions.

\section{Why some ambiguity does not disappear}
A long sequence of reports can remain uninformative about a parameter. In the rate example, if fixed and floating rates coincide throughout the observed period, additional interest totals may never distinguish the contract types. Sample size alone does not resolve an unidentified direction.

The set-identification analysis of \citet[Sections 2.1--2.3]{giacomini2021} separates parameters identified by the observable law from aspects of the prior that data cannot revise. The accounting application is conditional: one must first establish which observations depend on which hidden quantities. It is not enough to invoke a general consistency slogan and expect all uncertainty to vanish.

Persistent ambiguity can also become irrelevant to one query without being resolved. After final repayment, the old rate clause may no longer affect future cash. It may still matter for explaining historical interest, estimating management behavior or evaluating a counterfactual refinancing decision. Removing it from an active cash-flow state is different from asserting that the evidence identified it.

\section{Managers need not share the analyst's ignorance}
Public disclosures may leave $\theta$ unknown to an analyst even though management signed the contract. A management policy can legitimately depend on that private fact. If managers refinance only floating-rate debt after a rate shock, their actions differ across histories that the analyst cannot distinguish. That is not foresight; it is a specified information advantage.

Let $I_t^a$ denote actor $a$'s information history and let $\pi_t^a$ map it to a decision. Then
\begin{equation}
 a_t^a=\pi_t^a(I_t^a).
 \label{sim:actor-policy}
\end{equation}
Two simulated worlds with the same $I_t^a$ must receive the same action under a deterministic policy. For a randomized policy, they must receive the same conditional action distribution. The simulator's access to hidden state does not entitle the actor to use it.

Different actors can have different information. A lender may observe covenant certificates before public investors. A subsidiary may know its operating cash before a parent receives a report. A regulator may receive nonpublic prudential data. The model should specify these information channels where they affect behavior; assuming universal knowledge can make financing appear available too early.

\section{A simple foresight counterexample}
A manager must choose at date zero between two production plans. Plan A yields 10 if next year's demand is high and zero if low. Plan B yields 6 in either state. Assume high and low demand are equally likely and the manager learns demand only after choosing. The expected values are 5 and 6, so the expected-value policy selects B.

A simulator that first draws demand and then chooses the best plan in that realized world obtains 10 in the high state and 6 in the low state, averaging 8. It has valued perfect advance information. The arithmetic is correct for a different decision problem. The target problem permits an average of 6.

The discrepancy is not repaired by drawing more scenarios. It is caused by an information violation. A nonanticipativity check must require the same date-zero action across worlds sharing date-zero information. In a tree formulation this often appears as equality constraints on decisions at indistinguishable nodes. In a symbolic representation it can be enforced by restricting the policy's inputs.

\section{Signals can have economic value}
Suppose the manager can pay 1 at date zero for a perfectly informative demand signal before selecting A or B. With the signal the gross expected payoff is 8 and the net payoff is 7. Without it the payoff is 6. The signal is worth 2 gross and 1 net in this example. These calculations use a declared distribution and feasible action set.

The same logic applies to obtaining a debt schedule before committing liquidity, but real disclosures may be noisy, delayed or nonpublic. A signal's value depends on whether the actor can act after receiving it and whether the relevant action remains feasible. A maturity note arriving after repayment cannot improve that repayment decision, even if it improves a retrospective explanation.

Without probabilities, one can compare worst-case payoffs or regret before and after a partition of the family. The result is not an expected value of information. An interface should name the criterion because different criteria can recommend different research priorities.

\section{Simulated learning does not revise actual evidence}
A scenario may include a hypothetical disclosure at a future date. Inside that run, an actor updates its information and changes its action. The hypothetical report must not enter the analyst's actual evidence store as a fact. Otherwise running a scenario can contaminate later reconstructions with imagined information.

Separate identifiers should distinguish actual source assertions, simulated observations and analyst assumptions. A scenario observation can refer to the source model that generated it and to the hidden state of that run. It cannot acquire the authority of a filed report merely because its text looks like one.

The distinction also matters for conversational systems. An explanation generated from a selected witness should say that it is one compatible history. If the explanation is later ingested as evidence without its conditional status, a possible event becomes a purported fact. Provenance and type checks should prevent that feedback.

\section{Learning can change the appropriate compression}
Two histories may have equal current cash and equal next-period payments, yet differ in the chance of a revealing report. If the report affects a later refinancing choice, their future decision values can differ. A merge based only on current accounting totals discards this difference.

Suppose two contract states both pay 5 next year. In one state, a scheduled certificate will reveal a restrictive covenant before the next refinancing date; in the other, no such information arrives. If management or lenders condition their offers on that certificate, equal immediate cash does not imply equal future financing. The observation process belongs in the transition model used for compression.

This is why exact aggregation later requires equality of the joint law of successor state, observations and requested outputs for each allowed action. Preserving a marginal cash distribution is too weak. A graph may share current arithmetic while retaining separate information states for future decisions.



\section{Recursive ambiguity and persistent contracts}
Bayesian conditioning and a fresh worst-case choice at every date answer
different questions. Suppose a persistent binary type $\theta$ is either zero
or one, and tomorrow a fair binary signal $X$ unrelated to that type is observed.
The admitted laws are $Q_0$ and $Q_1$: under $Q_i$, $\theta=i$ with certainty
and $X$ remains fair. For the loss $Z=\ind\{X=\theta\}$, either law gives
expected loss $1/2$, so
\begin{equation}
\max_{Q\in\{Q_0,Q_1\}}\E_Q[Z]=\tfrac12,\qquad
\E_X\!\left[\max_{i\in\{0,1\}}\E_{Q_i}[Z\mid X]\right]=1.
\end{equation}
The second expression selects type zero after signal zero and type one after
signal one. It therefore evaluates a new law with $\theta=X$. That law is
absent even from the convex hull of the original two laws: every mixture
retains independence between the signal and type. No type has physically
changed along a realized path, yet dependence has changed across possible
paths. Preserving path support alone would miss the error.

Shapiro defines rectangularity through equality between static worst-case
evaluation and its nested conditional counterpart
\cite[Definition~2.1, Theorem~2.1 and proof, Section~5]{litcloseshapiro}.
His finite construction pastes conditional distributions across histories; it
can enlarge the family. A dynamic program for that enlarged family solves
a changed robustness question. Its convenience cannot justify claiming that
it solved the original persistent-type problem.

Terminology also requires care. Pichler and Shapiro call a product family a
rectangular setting, but explicitly show that its original worst-case
functional and its composite conditional functional can differ
\cite[Section~4.2, Example~4.7]{litclosepichlershapiro}. Computing a conditional
law-based risk measure can also differ from conditioning an unconditional
ambiguity set. Take six equiprobable states and an observed three-state
group with losses $0,1,2$. Conditional expected shortfall at level $1/2$
within that group is $5/3$. The unconditional six-state risk envelope permits
probabilities at most $1/3$ per state. It can assign $1/6$ to the loss-2 state,
zero to the other two group members, and $5/18$ to each state outside the
group. This is an admissible probability law; after conditioning on the group
it assigns probability one to loss 2. The conditional worst case over that
envelope is therefore 2. The difference is between two precise conditional
risk definitions, not an arithmetic error
\cite[Example~3.2 and equations~(3.1)--(3.5)]{litclosepichlershapiro}.

Policy consistency and equality of risk values are distinct as well. Xin and
Goldberg give moment-constrained inventory examples in which an initially
optimal policy ceases to be optimal after reoptimization, and another in which
the optimal policy sets agree although the two optimal values differ
\cite[Definitions~3.1--3.3, Observations~5--7, proofs in Sections~6.3--6.5]{litclosexingoldberg}.
The company model must specify whether management can commit to a policy,
what information each later decision uses, and which family of joint laws is
retained. For a finite reference problem, enumerate complete admissible laws
and observation-measurable policies before choosing an approximation. For
larger problems, a set of beliefs over persistent types or an explicit
precommitment formulation can retain those restrictions, but its recursive
sufficiency and computational cost need a separate derivation. Repeated
worst-case optimization remains useful as a stress test when the changed
family is stated explicitly.



\section{Persistent parameters and particle approximation}
A particle approximation can reduce computation after a probability model is
chosen. It does not create that model. A type with probability $0.01$ has
probability $0.99^{100}>1/3$ of disappearing from 100 independent multinomial
resampling draws. No new evidence is involved in this loss. Rare but material
types therefore need explicit support checks, stratification or separate
conditional computations, with the resulting approximation stated.

For state-space models with persistent unknown parameters, SMC$^2$ combines
parameter particles and state filters, with particle-MCMC rejuvenation
\cite[Section~3.2, equation~(7), Proposition~1, Appendix~A]{litclosesmc2}. Its
extended target has the desired parameter posterior as a marginal under the
source assumptions. Finite particles still approximate that posterior; exactness
of a target marginal does not mean that one finite run represents every
possible type. The unbiased full particle likelihood estimate used in this
construction must be distinguished from an arbitrary estimated likelihood
ratio.

The inspected author-linked modern implementation creates a state filter for
each parameter particle, runs that filter through the current observation date
when evaluating a proposed parameter, and incorporates a likelihood correction
when increasing the state-particle count. These operations have real time and
memory costs. Retaining only filtering arrays differs from retaining complete
genealogies. This is a candidate for explicitly probabilistic filtering, not
a replacement for symbolic unweighted historical ambiguity.


\section{Changing particles without changing the target}
\label{sim:inference-exchange}

Increasing the number of particles can change an inference result even before
new evidence arrives. The reason is that a likelihood estimate and its random
construction belong to the state used by particle MCMC. Replacing the estimate
with another unbiased estimate does not automatically preserve their joint
distribution. A test that examines only current company-type probabilities
can miss the error.

Fix an evidence revision $v$ and a persistent type $\theta$. Write its prior
as $p(\theta)$ and likelihood as $L_v(\theta)$. A particle representation of
size $N$ uses random variables $u$ with law $m_N(u\mid\theta)$ and produces
a nonnegative unbiased estimate $Z_{N,v}(\theta,u)$ of that likelihood.
Assume the normalizing constant $C_v$ is finite and positive. The joint target
weights the random construction by the estimate it produces:
\begin{equation}
 \widetilde\pi_{N,v}(\theta,u)=
 \frac{p(\theta)m_N(u\mid\theta)Z_{N,v}(\theta,u)}{C_v}.
 \label{sim:extended-particle-target}
\end{equation}
Integrating over $u$ replaces $Z_{N,v}$ by its expectation $L_v$, giving the
desired type posterior. This is the extended-target argument in the retained
SMC$^2$ source, arXiv 1101.1528v3, 27 January 2012
\citep[equation (7), Proposition 1 and Appendix A]{litclosesmc2}.
Unbiasedness establishes this marginal relation; it does not make an
arbitrary finite particle population equal to the posterior.

Now change representation from $N$ to $M$ and draw fresh randomness $u'$ from
$m_M(\cdot\mid\theta)$. The augmented proposal law is
$\widetilde\pi_{N,v}(\theta,u)m_M(u'\mid\theta)$. With backward law
$m_N(u\mid\theta)$, the desired augmented target is
$\widetilde\pi_{M,v}(\theta,u')m_N(u\mid\theta)$. Their ratio is
\begin{align}
 w(\theta,u,u')
 &=\frac{\widetilde\pi_{M,v}(\theta,u')m_N(u\mid\theta)}
         {\widetilde\pi_{N,v}(\theta,u)m_M(u'\mid\theta)}\\
 &=\frac{Z_{M,v}(\theta,u')}{Z_{N,v}(\theta,u)}.
 \label{sim:exchange-weight}
\end{align}
The prior, normalizer and auxiliary proposal laws cancel because the type and
evidence revision are unchanged. This recovers the exchange correction in
Section 3.6.1 of the inspected source. The ratio requires the relevant support
conditions. The following example uses strictly positive estimates, so it
does not hide a zero-denominator problem inside that cancellation.

Take two types with prior probability $1/2$ each and true likelihood one.
For type zero the old estimate is $1/2$ or $3/2$ with equal probabilities;
the new estimate is $1/4$ or $7/4$, again equally likely. Type one's estimate
is always one. Each estimator is unbiased. Order the new joint states as
$(0,1/4),(0,7/4),(1,1)$. The likelihood-weighted target in
equation~\eqref{sim:extended-particle-target} and the distribution produced by
unweighted replacement are respectively
\begin{equation}
 \pi=\left(\frac1{16},\frac7{16},\frac12\right),\qquad
 q=\left(\frac14,\frac14,\frac12\right).
\end{equation}
Both assign probability $1/2$ to type one. Immediate marginal agreement
therefore passes a test that the joint target fails.

To see the consequence, make one valid particle marginal Metropolis--Hastings
move. Propose the other type, draw its new estimate, and accept with
$\min(1,Z_{\rm proposed}/Z_{\rm current})$. From $(0,1/4)$ the proposal is
accepted certainly. From $(0,7/4)$ it is accepted with probability $4/7$.
From type one, the low-estimate proposal has total accepted probability
$(1/2)(1/4)=1/8$, and the high proposal has probability $1/2$. The transition
matrix, with the same state order, is thus
\begin{equation}
 K=\begin{pmatrix}
 0&0&1\\
 0&3/7&4/7\\
 1/8&1/2&3/8
 \end{pmatrix}.
\end{equation}
The correct joint law satisfies $\pi K=\pi$ and detailed balance. Starting
instead from $q$, the probability of type one after the move is
\begin{equation}
 (qK)_3=\frac14+\frac14\frac47+\frac12\frac38
       =\frac{65}{112}
       =\frac12+\frac9{112}.
\end{equation}
The transition is valid for its stated target; the wrongly initialized joint
distribution is the problem. This is an exact one-step counterexample, not a
claim that a valid PMMH kernel has the wrong stationary distribution. Applying
the exchange weight restores the desired joint law in the finite enumeration.

The inspected author-linked implementation includes the exchange
correction \citep[\texttt{SMC2.exchange\_step}, pinned source]{particlesgap2026}.
Its new filter is run through the preceding observation before
the next incremental weight is added; processing that next observation in
both places would count it twice. Rejected Metropolis proposals retain their
existing filter state, and repeated resampled selections receive distinct
mutable filter copies. These are scoped source-code correspondences, not a
run of the full modern package or a performance comparison. The modern
waste-free extension requires its own dependence and finite-run analysis.

\section{A correction can restore an excluded history}
Changing the particle representation preserves the likelihood being computed.
Correcting a historical disclosure can change that likelihood itself. Suppose
two histories have prior probabilities $(1/2,1/2)$. An old hard observation
has likelihoods $(1,0)$, producing posterior $(1,0)$. Its replacement has
likelihoods $(1/4,3/4)$. Recomputing from the prior gives
\begin{equation}
 \pi_{\rm replacement}=(1/4,3/4).
\end{equation}
Multiplying the restricted old posterior by the replacement likelihood and
normalizing instead gives $(1,0)$. The operation conjoins the old observation
with the new one. It cannot recover the history the old assertion excluded.
The new-to-old likelihood ratio is also undefined on that restored history,
where the denominator was zero.

A replacement therefore requires replay from a sufficiently informative
preserved checkpoint, or a separately justified bridge with the required
support. The inference request must identify the evidence revision and which
assertion is replaced. Merely attaching a corrected document to an existing
particle population does not establish the corrected posterior. This is the
probabilistic counterpart of retaining alternative histories that become
admissible after a source correction.

The source's quantitative guarantees require additional premises. Its
Appendix B uses uniform transition and observation-density ratio bounds and
posterior regularity conditions, and its refreshed-particle idealization is
stronger than a finite PMCMC move
\citep[Proposition 2 and Appendix B]{litclosesmc2}. A permanent deterministic
type, a hard zero-likelihood accounting observation or an unidentified
parameter direction can violate the intended transfer assumptions. Exact
finite target calculations establish the distinctions above; they do not
establish company-model mixing, scaling, likelihood identification or
forecast accuracy.


\section{Replaying a declared evidence vintage}
\label{sim:finite-vintage}
The preceding correction becomes an executable calculation once the inference
request says exactly which information remains active. A document identifier
alone is insufficient. Two documents may repeat one observation, and a new
document may replace an earlier assertion rather than supply a second
independent measurement. The model must retain these relations before it
multiplies likelihoods.

Consider a finite set of persistent company types $h\in H$, with declared prior
masses $p(h)$. Each active observation occupies a named slot and supplies a
likelihood $L_j(h)$. We assume that these observation blocks are conditionally
independent given $h$ in this example. Under that assumption, for a selected
evidence vintage $v$ with active slots $A_v$,
\begin{equation}
 Z_v=\sum_{h\in H}p(h)\prod_{j\in A_v}L_j(h),\qquad
 \pi_v(h)=\frac{p(h)\prod_{j\in A_v}L_j(h)}{Z_v}.
 \label{sim:finite-vintage-posterior}
\end{equation}
The normalizer must be positive. A zero value means that this finite model
assigns zero likelihood to the selected evidence; it does not justify returning
an arbitrary distribution. Missing prior probabilities or likelihoods instead
leave the calculation unresolved. Neither circumstance supplies probabilities
for a family that was specified only as a feasible set.

The selected reference, \texttt{FiniteVintageInference/1}, retains the named
types, exact rational prior, observation-model identity and factorization
assumption. A request also names the actor and information cutoff, the active
observation versions and their replacement ancestry. Each observation retains
its source content, information identity, slot, likelihood vector, supplied
actor-availability time and any superseded version. Its content identifier
changes when any of those fields changes. The semantic information identity
answers a different question: whether two records purport to describe the
same informational contribution. A replacement keeps that identity even when
its likelihood or source changes.

The supplied availability time must be no later than the request's cutoff.
An unknown time remains unresolved; a later time is ineligible for that
request. The reference uses the existing exact fixed-offset timestamp
conversion, including fractional seconds, rather than rounding time to a
calendar day. Its convention excludes leap seconds and unspecified offsets.
These checks establish consistency with declared availability. Establishing
when an actual investor, manager or system could obtain the information still
requires a source-selection procedure.

Replacement selects a new head for an existing slot and preserves the chain
to the earlier record. Retraction removes that slot's likelihood contribution,
equivalently replacing it by one for every type. The request cannot select two
heads for one slot, use a missing predecessor, or treat a different information
identity as a revision of the same observation. Several unresolved candidate
revisions therefore require several conditional requests. Selecting a head
does not establish that it is the authoritative version.

The numerical consequences are immediate in the two-type example. After the
correction, the posterior is $(1/4,3/4)$. Suppose a separately declared,
conditionally independent signal has likelihoods $(1/2,1/4)$. Multiplying
gives unnormalized masses $(1/8,3/16)$, whose sum is $5/16$, and hence
\begin{equation}
 \pi_{\rm corrected+signal}=(2/5,3/5).
\end{equation}
Retracting the corrected observation leaves only the new signal; with the
original equal prior the resulting posterior is $(2/3,1/3)$. Retracting the
only observation restores the prior. A type assigned zero prior mass remains
at zero unless the prior itself is changed explicitly. Restoring support lost
through an erroneous observation is distinct from inventing support absent
from the declared model.

Repackaging the corrected observation under another slot would produce a
different answer. With the equal prior, counting its likelihood twice gives
masses proportional to $(1/16,9/16)$ and posterior $(1/10,9/10)$. The reference
rejects two active observations that declare the same semantic information
identity. An earlier implementation checked distinct slots but omitted this
condition: its arithmetic tests passed while the declared information
contract remained incomplete. The repaired example makes the resulting
overconfidence visible. False distinct identities can still conceal repeated
information, just as a false independence assumption can make a well-formed
calculation substantively wrong.

A checkpoint retains its complete request and computed result. Reuse requires
the same model and prior, actor and exact cutoff, with the old active records
forming an unchanged prefix of the new sequence. New independent records can
then update that prefix. A replacement that changes the prefix instead causes
replay from the preserved prior. This avoids the failed operation that starts
from $(1,0)$ and attempts to recover the excluded type by reweighting it.
Corrupt checkpoint contents or incompatible model and information definitions
are rejected rather than silently relabelled.

In this finite reference the old checkpoint calculation is itself recomputed
before reuse, so the test establishes the posterior relation without claiming
a computational saving. Safe reuse from a larger store would require a proved
sufficient dependency set. Particle checkpoints would additionally retain the
likelihood estimate, proposal and ancestor state required by the extended
target. A vector of posterior weights alone is insufficient.

The reference compares complete masses and normalizers over a finite rational
test grid and explicit contrary cases, including a correction, retraction,
duplicate information, unknown times and a later fractional-second timestamp.
An uninformative observation leaves both types alive. This illustrates a
representation in which uncertainty can persist without enumerating a new
copy of every historical node. It establishes the exact finite calculation
under its declared premises; continuous inference, source completeness and
the suitability of the likelihood for a real company remain further questions.


\section{A correction reaches different decision makers at different times}
\label{sim:actor-vintage-decision}
Replaying the corrected history answers an inference question. It does not
yet say which decision maker could use the correction. Suppose the analyst
learns on 1 February 2026 that the earlier observation was wrong, while the
manager receives that correction on 1 March. On 15 February they possess
different information even though they are studying the same company. A
simulation of the manager's choice must use the manager's information.

Continue the hypothetical fixed-versus-floating example. Under the original
observation the posterior is $(1,0)$; under the correction it is
$(1/4,3/4)$. Suppose an unhedged exposure costs zero for the fixed type and
20 for the floating type. A perfect hedge costs 6 in either case. These are
supplied losses in the same units, and both actions are assumed feasible.
They are enough to define this comparison without inventing an observed
hedging price or an accounting transaction.

Let $p$ denote the probability of the floating type under the particular
actor's available information. The expected losses are therefore
\begin{equation}
  L(\text{unhedged}\mid p)=20p,\qquad
  L(\text{hedged}\mid p)=6.
  \label{sim:actor-hedge-loss}
\end{equation}
The original posterior gives an unhedged loss of zero. The corrected
posterior gives 15, so the hedge has the smaller declared expected loss.
At $p=3/10$ the two losses both equal 6 and both actions remain minimizers.
An implementation must retain that tie unless a separate choice rule is
supplied. The manager on 15 February still uses the old posterior, while
the analyst can already select the hedge under the corrected one.

The selected information procedure takes an actor and an exact cutoff,
selects the available unreplaced observation in each information slot, and
passes that request to the finite evaluator. Adding the future correction
does not change the manager's stored 15 January request or its decision.
Later knowledge can motivate a new decision, but it must not silently
replace the information used to reproduce an earlier one.

Unknown availability cannot be treated as known absence. Two incomparable
available revisions also leave the request unresolved; choosing the newer
file arbitrarily would add an authority rule. The checked examples retain
these distinctions, refuse corrupt observation content and fabricated
posterior results, and leave unknown action losses unresolved. They select
actions from the declared feasible set; they do not execute a hedge or infer
its causal effect. Actual actor access, registry completeness, instrument
terms and feasible financing still need evidence.




\section{A proposed information service}
AccDSL should represent evidence views and actor information views separately. An evidence view selects actual assertions by availability, revision and scope. An actor view specifies observable state variables, private contractual facts and signal channels at each decision date. Policies receive only that view and permitted memory.

Forward execution first applies events available at the current boundary, constructs observations, updates the relevant information state, and then permits decisions according to the stated timing. Alternative orderings are possible, but they must be explicit and tested. Moving a decision across an observation boundary changes the economic problem.

The service should refuse ambiguous timing where it affects the answer. It can still return a conditional comparison of both orderings. Such a comparison is preferable to silently selecting whichever ordering makes the simulation easier. Exact timestamps, time zones and publication lags matter when events occur on the same calendar date.

\section{Tests that discriminate information errors}
A useful test duplicates a world, changes an unobserved future shock and confirms that the current action remains unchanged. Another changes a privately known contract term and confirms that an informed manager's policy may change while a public analyst's evidence view does not. A third inserts a later restatement and checks that an earlier frozen forecast remains reproducible.

These are semantic tests. They do not require a fitted behavioral model. They establish that information restrictions are implemented as specified. Whether a particular manager actually follows the chosen policy is a separate empirical question. Keeping that distinction prevents a passing information test from being presented as behavioral validation.

A final test removes all informative observations and confirms that persistent uncertainty remains. A system that gradually collapses to a single history merely because a particle filter resampled or an optimizer selected one witness has confused computational concentration with evidence. The retained family and the approximation used for calculation must remain distinguishable.
